Eligibility depends on the contract and the insurer's expectations at inception.
The policyholder must participate in a clearly identified pool, with the
required substantial share of fair-value returns and substantial variability
in payments. The variability assessment covers the contract's duration on a
probability-weighted present-value basis. A minimum guarantee can bind in some
scenarios without automatically excluding direct participation. Later market
movements do not reopen the classification; contract modification follows its
own rules. The identified reference pool also need not be held by the insurer,
so the contractual link and the actual asset holdings require separate records
\citep[paragraphs B101--B108 and 72]{aasb17}.

\clearpage
A small settlement example shows why the guarantee matters even when the
stated share is fixed. Consider only a maturity-benefit component, in one
currency at one date. Assume a held pool allocated to this component has amount
$U$, and the policyholder is owed the larger of 90 and $0.9U$. These are
illustrative contractual choices. The amount left from the pool after payment
is $U-\max(90,0.9U)$. Writing the payment as
$0.9U+\max(90-0.9U,0)$ shows that the amount left is
$0.1U-\max(90-0.9U,0)$: the nominal retained share less the extra payment
needed to meet the minimum.

\begin{center}
\begin{tabular}{@{}rrr@{}}
\toprule
Pool amount & Policyholder payment & Amount left\\
\midrule
80 & 90 & $-10$\\
100 & 90 & 10\\
120 & 108 & 12\\
\bottomrule
\end{tabular}
\end{center}

Moving from 80 to 100 leaves 20 more after payment: the ordinary share rises
by 2 and the extra floor payment falls by 18. Moving from 100 to 120 leaves
only 2 more because the minimum no longer binds. The negative amount at 80
is a funding shortfall for this component. These unweighted settlement
outcomes explain the exposure; they do not establish insurance-contract
eligibility, a present value or the insurer's profit. Other benefits and costs
remain outside this component example, and its held-pool assumption does not
replace the standard's separate accounting for investment assets
\citep[paragraphs B66(a), B107--B108]{aasb17}.

The accounting calculation therefore needs more than this payoff schedule.
Paragraph B104 defines the variable fee as the entity's share of underlying
fair value less fulfilment cash flows that do not vary based on underlying
returns. Changes in the obligation to pass through the underlying fair value
do not adjust CSM, whereas the entity-share changes relate to future service.
Financial-guarantee changes have their own future-service treatment under
B113(b). Paragraph 45 retains the loss-component and risk-mitigation exceptions,
as well as new-contract, currency and service-allocation movements. Applying
those rules requires the group's measured cash flows, discount and risk
adjustments, remaining CSM and loss component, service information and any
eligible mitigation designation. The original increase of 2 remains a valid
share-component calculation; the complete CSM movement needs these additional
inputs \citep[paragraphs 32--38, 45, 48--50 and B104--B118]{aasb17}.
